\section{Related Work}
\label{sec:related}

\paragraph{Vision-language models for UI understanding.}
The application of vision-language models to graphical user interfaces has accelerated rapidly. Early work focused on OCR-centric approaches that extracted text from screenshots before reasoning over the extracted tokens~\cite{tang2023unifying}. More recent systems operate directly on pixel inputs: ScreenAI~\cite{baechler2024screenai} trains a unified model for screen annotation and question answering by leveraging automatically generated training data across multiple screen types. Ferret-UI~\cite{you2024ferretui} extends the Ferret architecture~\cite{you2023ferret} with adaptive gridding to handle the high-resolution, aspect-ratio-diverse images typical of mobile interfaces, enabling fine-grained referring and grounding on UI elements. CogAgent~\cite{hong2024cogagent} introduces a cross-attention mechanism between a low-resolution general encoder and a high-resolution screen encoder, achieving strong results on both web and mobile UI benchmarks. These models demonstrate impressive capabilities on standard UI understanding tasks, but they are invariably evaluated on canonically rendered screenshots. Our work complements these advances by providing a diagnostic benchmark that measures how robust these capabilities are when the visual presentation of the same underlying content varies.

\paragraph{UI and document understanding benchmarks.}
Several benchmarks target visual question answering over screenshots and documents. ScreenQA~\cite{hsiao2022screenqa} provides question-answer pairs grounded in mobile application screenshots, while WebSRC~\cite{chen2021websrc} offers a structural reading comprehension benchmark for web pages that includes HTML structure alongside screenshots. VisualWebBench~\cite{liu2024visualwebbench} proposes a multimodal benchmark for evaluating VLMs on web page understanding across dimensions including OCR, element grounding, and web-specific tasks. DocVQA~\cite{mathew2021docvqa} and InfographicVQA~\cite{mathew2022infographicvqa} target document and infographic comprehension, respectively. While these benchmarks comprehensively evaluate model accuracy on diverse UI and document types, none of them systematically vary the visual presentation of the underlying content to assess robustness. \bench{} fills this gap by introducing a controlled drift taxonomy with calibrated severity levels, enabling researchers to measure not just whether a model can answer correctly, but whether it can do so consistently across visual variations.

\paragraph{Visual robustness benchmarks.}
The study of model robustness to visual perturbations has a rich history in image classification. ImageNet-C~\cite{hendrycks2019imagenetc} introduced 15 corruption types (noise, blur, weather, digital) at five severity levels, establishing a standard protocol for evaluating classification robustness. The companion CIFAR-10-C benchmark~\cite{hendrycks2019cifarc} extended this protocol to a smaller-scale setting. Subsequent work has proposed robustness benchmarks for object detection~\cite{michaelis2019objectdetection}, semantic segmentation, and other vision tasks, consistently finding that models trained on clean data degrade substantially under distribution shift. In the document domain, Larson \etal~\cite{larson2022docrobust} studied the effect of image degradation on OCR systems, while Borchmann \etal~\cite{borchmann2021due} evaluated document understanding models under varied formatting. Our work extends this robustness evaluation paradigm to the UI understanding setting, replacing pixel-level corruptions with semantically meaningful layout and style perturbations that mirror real-world UI variability.

\paragraph{Grounded visual question answering.}
Grounded VQA requires models to not only answer questions but also localize the evidence supporting their answers. Zhu \etal~\cite{zhu2016cvpr} introduced grounded VQA as a task requiring spatial evidence for answers in natural images, while later work extended grounding to document and chart understanding~\cite{masry2022chartqa,chen2023palix}. In the UI domain, grounding takes on particular importance because the visual evidence region provides a verifiable audit trail: if an agent extracts a revenue figure from a dashboard, a human reviewer needs to know which chart or table cell the figure came from. Despite its importance, grounding quality is rarely measured in UI benchmarks. \bench{} includes ground-truth bounding boxes for every QA pair, enabling evaluation of evidence localization as a first-class metric alongside answer accuracy.
